\documentclass[11pt]{article}
\usepackage{problem_set_style}
\usepackage{kotex}
\usepackage{booktabs}
\usepackage{array}
\begin{document}
\heading{}{샘플}{}{Spring 2027}{프로젝트 2 · Abstract 샘플 (11주차 제출, 1쪽)}

\noindent{\small 아래는 \textbf{가상의 팀}이 쓴 샘플입니다. 수치는 형식을 보이기 위한 예시이며 실제 실험 결과가 아닙니다. 형식과 분량의 기준으로만 참고하세요.}

\noindent\textbf{제목}: Fashion-MNIST에서 작은 CNN과 GBDT 비교: 딥러닝이 필요한가

\vspace{4pt}
\noindent\textbf{문제 정의.} 행 하나는 \textbf{옷 사진 한 장}(28$\times$28 흑백), 라벨은 10종류 중 하나이다. 클래스가 균형이라 주 지표는 정확도와 macro-F1이다.

\noindent\textbf{왜 딥러닝인가 (초안).} 픽셀 784개를 그대로 표처럼 넣으면 이웃 픽셀의 구조를 모델이 모른다. 합성곱이 이 구조를 쓰는지가 가설이며, GBDT와의 차이로 검증한다.

\noindent\textbf{데이터와 분할.} Fashion-MNIST 70{,}000장. 공식 train 60{,}000장을 50{,}000/10{,}000으로 나눠 train/validation으로 쓰고, 공식 test 10{,}000장은 마지막에 한 번만 쓴다.

\noindent\textbf{베이스라인.} 픽셀 784개를 입력으로 한 \textbf{GBDT}. 같은 train/validation으로 튜닝한다.

\noindent\textbf{계획.} 합성곱 2층의 작은 CNN을 만들고, 에포크 곡선으로 조기 종료 지점을 정한다. 데이터 증강(좌우 반전)과 드롭아웃의 효과를 구성 요소 제거 실험으로 확인한다.

\end{document}
